\section{Benchmark suite}
\label{app:suite}
\paragraph{Procedural fixtures (13).} cube (8\,$\times$\,8\,$\times$\,4 grid), slab, tower, L shape, T shape, wedge
(simple solids, expected valid with coverage at least 0.8); pyramid, sphere (overhangs); cylinder (budget-limited,
coverage at least 0.5); bridge; flat sheet (cannot be covered, expected partial); thin pole (zero valid builds,
expected no valid candidate); open box (non-watertight, expected rejected at parse time). The last is the one case
excluded from every mean; the thin pole is attempted and scores 0 for every generator.

\paragraph{Thingi10K subset (40).} Drawn with seed 0 from the Hugging Face mirror of Thingi10K among models that are
CC BY, CC BY-SA or public domain, closed, manifold, single-component solids with 100 to 200k faces, one file per
Thingiverse object, all with $+z$ up. Licence, author, name and source link are stored per model in the benchmark
records. The meshes are downloaded by a script and never redistributed.

\section{Hyperparameters and cost}
\label{app:hparams}
\begin{table}[h]
  \centering\small
  \begin{tabular}{@{}p{0.14\linewidth}p{0.84\linewidth}@{}}
    \toprule
    \textbf{Policy} & 4\,$\times$ Conv3d($3^3$, 32 ch) + GELU; candidate MLP 64-64-1 on pooled features + 13 (15 in RL)
    candidate features; STOP and value MLPs 64-1 on global mean + max; 97{,}058 parameters (101{,}283 with the RL
    layout and value head) \\
    \textbf{SFT} & teacher forcing, cross-entropy over candidates + STOP; AdamW $10^{-3}$, cosine schedule, batch 32,
    20{,}000 steps, grad clip; 8-fold symmetry augmentation; streaming item pool; free-run eval every 1000 steps on 64
    examples; v2 draws half of each pool refresh from real meshes and drops examples whose teacher covers under half
    the target \\
    \textbf{RL} & PPO, clip 0.2, entropy 0.01, KL 0.05 to the frozen initial policy, value coefficient 0.5, Adam
    $10^{-4}$, grad clip 1.0, $\gamma = 1$ (terminal reward), advantages normalised per batch; 250 iterations
    $\times$ 128 episodes, temperature 1, 8 forked CPU rollout workers, 3 epochs, minibatch 128; reward IoU after trim
    + 0.1 $\cdot$ stable; eval every 10 iterations on 64 examples (half real in v2); best checkpoint by eval IoU \\
    \textbf{Data} & v1: 2000 procedural + 150 real (at most 50k faces, teacher coverage at least 0.5), split
    1750 / 194 / 206. v2: 2000 procedural + 640 real (at most 200k faces, weak-teacher targets kept with their
    coverage recorded), split 2132 / 248 / 260 (test: 188 procedural, 72 real of which 30 weak-teacher) \\
    \textbf{Hardware} & one Apple M5 Pro laptop, 24\,GB; SFT and PPO updates on Metal, rollouts and all benchmark
    decoding on CPU for bit-reproducibility \\
    \textbf{Wall clock} & SFT 23 to 28 min per run; RL v1 2\,h\,25\,min, RL v2 1\,h\,35\,min (seed 0) and
    2\,h\,22\,min (seed 1), ablations about 2\,h\,25\,min each; the frozen benchmark 60\,s for the search and a few
    minutes per learned generator; the audit regenerates every record in a comparable time \\
    \bottomrule
  \end{tabular}
\end{table}

\section{The decode decision}
\label{app:stop}
\begin{table}[h]
  \centering\small
  \caption{The decode decision. Mean IoU after trim on the SFT test split (260 examples at the 100-piece cap: 188
  procedural, 72 real meshes, of which 30 have a teacher build covering under half the target), for the learned stop
  and for masked stop. Made before any benchmark run of the checkpoint.}
  \label{tab:stop}
  \footnotesize
  \begin{tabular}{llcccc}
    \toprule
    checkpoint & STOP & all & procedural & real & real, weak teacher \\
    \midrule
    greedy teacher & -- & 0.843 & 0.943 & 0.584 & -- \\
    SFT v1 & learned / masked & 0.770 / 0.775 & 0.863 / 0.864 & 0.527 / 0.543 & 0.175 / 0.184 \\
    SFT v2 & learned / masked & 0.766 / 0.782 & 0.870 / 0.876 & 0.493 / 0.535 & 0.123 / 0.179 \\
    RL v1 & learned / masked & 0.824 / 0.835 & 0.936 / 0.936 & 0.531 / 0.569 & 0.190 / 0.206 \\
    RL v2, seed 0 & learned / masked & 0.828 / 0.840 & 0.937 / 0.939 & 0.542 / 0.581 & 0.147 / 0.208 \\
    RL v2, seed 1 & learned / masked & 0.821 / 0.835 & 0.931 / 0.935 & 0.533 / 0.574 & 0.141 / 0.204 \\
    RL v2, no stability reward & learned / masked & 0.834 / 0.842 & 0.940 / 0.941 & 0.559 / 0.587 & 0.188 / 0.223 \\
    RL v2, no connectivity inputs & learned / masked & 0.766 / 0.805 & 0.885 / 0.906 & 0.455 / 0.543 & 0.129 / 0.197 \\
    \bottomrule
  \end{tabular}
\end{table}

\section{Stability model}
\label{app:stability}
Bricks are rigid and loaded by their own weight (2.32\,g per 2\,$\times$\,4 brick, scaled by footprint area).
Each shared stud cell between an upper and a lower brick carries vertical force at four sub-points near the cell
corners; ground cells carry compression only. Per brick, vertical force balance and torque balance about $x$ and $y$
are imposed. A linear program finds the smallest tension utilisation $t$ such that every joint force lies in
$[-t\,T/4,\,+\infty)$ with $T = 0.98$\,N per stud cell. StableLego~\citep{liu2024stablelego} measured 0.98\,N as the
capacity of one contact point (three to four per stud), so dividing it over four points per stud makes this model
about four times more conservative than theirs; we keep the conservative reading because no capacity was measured in
this project.
$t \le 1$ is reported as stable under the model; an infeasible program even with unlimited tension means the centre
of mass has left the support (topple). The model has no lateral loads, shear, friction, deformation, hand forces or
dynamics, and its constants were not measured on any physical brick in this project. One program takes about 20\,ms
for a 75-piece build.

\section{Assembly planner}
\label{app:planner}
Rules: one brick per step, inserted straight down onto a horizontal work surface, starting from a ground brick. A
brick may be added when it does not collide, attaches to the assembled component (ground contact or studs of a placed
brick directly below), its vertical insertion path is clear, and it does not seal an unplaced brick beneath it. The
last rule is an exact dead-end test, which makes greedy selection complete: if any full order exists, the first
unplaced brick of that order is always attachable and safe, so the planner never backtracks. Ties break on lowest
layer, then most placed supporters, then id, which gives a layer-by-layer order. For this catalogue (same-height
upright bricks, every non-ground brick supported from below) every structurally valid build has a planned order; the
validator replays the steps and checks the insertion path.

\section{Reproducibility}
\label{app:repro}
All code, benchmark records, audit logs and figure scripts are in the repository. The commands below regenerate the
paper's numbers from a fresh checkout (training needs the Thingi10K download, about 1.5\,GB, and the laptop GPU):
\begin{verbatim}
make bench                                   # frozen greedy baseline + audit
make sft-train sft-bench rl-train rl-bench   # v1 ladder
make data-v2 sft-train-v2 rl-train-v2 \
     rl-bench-v2                             # v2 ladder, seed 0
make v2-seed1 rl-bench-v2-seeds              # v2 ladder, seed 1
make rl-ablations rl-bench-v2-ablations      # ablations
make paper-builds paper-figs paper-check \
     paper-pdf                               # figures, macros, checks, this PDF
\end{verbatim}
Every number in the main tables is a macro written by the figure script from the benchmark records, so the text
cannot drift from the data. Decoding runs on CPU and is bit-identical across processes; voxelization pins the random
state used by the containment test. The benchmark files are the six \code{.jsonl} files under \code{docs/benchmarks/} named in Section~\ref{sec:benchmark}.

\section{Per-case results}
\label{app:percase}
% generated by paper/scripts/figs.py tables; do not edit
Per-case IoU after trim at the 200-piece cap (0 for a missing or invalid build). RL columns use the masked stop; cells is the number of occupied target cells out of 2048 (512 for the cube).
\begin{center}\footnotesize\setlength{\tabcolsep}{4pt}
\begin{tabular}{llrcccccc}
\toprule
case & source & cells & search & SFT v2 & RL v2 s0 & RL v2 s1 & no stab. & no conn. \\
\midrule
bridge & procedural & 224 & 0.964 & 0.964 & 0.964 & 0.964 & 0.964 & 0.964 \\
cube & procedural & 256 & 1.000 & 0.641 & 1.000 & 1.000 & 1.000 & 1.000 \\
cylinder & procedural & 608 & 1.000 & 1.000 & 1.000 & 1.000 & 1.000 & 1.000 \\
flat\_sheet\_impossible & procedural & 160 & 0.050 & 0.050 & 0.050 & 0.050 & 0.050 & 0.050 \\
l\_shape & procedural & 540 & 1.000 & 1.000 & 1.000 & 1.000 & 1.000 & 1.000 \\
pyramid & procedural & 336 & 0.946 & 0.911 & 0.946 & 0.928 & 0.918 & 0.911 \\
slab & procedural & 480 & 1.000 & 1.000 & 1.000 & 1.000 & 1.000 & 1.000 \\
sphere & procedural & 392 & 0.975 & 0.959 & 0.957 & 0.954 & 0.959 & 0.964 \\
t\_shape & procedural & 288 & 1.000 & 1.000 & 1.000 & 1.000 & 1.000 & 1.000 \\
thin\_pole & procedural & 0 & 0.000 & 0.000 & 0.000 & 0.000 & 0.000 & 0.000 \\
tower & procedural & 32 & 1.000 & 1.000 & 1.000 & 1.000 & 1.000 & 1.000 \\
wedge & procedural & 192 & 1.000 & 0.990 & 0.995 & 0.995 & 0.990 & 0.995 \\
T100070 & thingi10k & 1070 & 1.000 & 0.938 & 1.000 & 1.000 & 0.999 & 0.998 \\
T100894 & thingi10k & 84 & 0.913 & 0.905 & 0.907 & 0.905 & 0.907 & 0.907 \\
T105120 & thingi10k & 56 & 0.200 & 0.267 & 0.194 & 0.271 & 0.194 & 0.138 \\
T108578 & thingi10k & 68 & 0.057 & 0.000 & 0.057 & 0.057 & 0.057 & 0.057 \\
T108884 & thingi10k & 276 & 0.953 & 0.931 & 0.933 & 0.953 & 0.971 & 0.932 \\
T110201 & thingi10k & 32 & 0.250 & 0.250 & 0.250 & 0.273 & 0.250 & 0.250 \\
T111878 & thingi10k & 60 & 0.109 & 0.109 & 0.109 & 0.109 & 0.109 & 0.108 \\
T1130077 & thingi10k & 205 & 0.698 & 0.718 & 0.718 & 0.718 & 0.718 & 0.718 \\
T113134 & thingi10k & 108 & 0.404 & 0.393 & 0.460 & 0.327 & 0.414 & 0.447 \\
T130960 & thingi10k & 184 & 0.967 & 0.968 & 0.968 & 0.947 & 0.957 & 0.968 \\
T1335002 & thingi10k & 400 & 0.980 & 0.946 & 0.950 & 0.947 & 0.941 & 0.948 \\
T1417953 & thingi10k & 256 & 1.000 & 1.000 & 1.000 & 1.000 & 1.000 & 1.000 \\
T144106 & thingi10k & 64 & 0.300 & 0.210 & 0.308 & 0.316 & 0.316 & 0.333 \\
T1514900 & thingi10k & 92 & 0.706 & 0.593 & 0.652 & 0.702 & 0.652 & 0.626 \\
T200212 & thingi10k & 365 & 0.989 & 0.981 & 0.992 & 0.986 & 0.984 & 0.981 \\
T225935 & thingi10k & 290 & 0.907 & 0.852 & 0.883 & 0.865 & 0.888 & 0.837 \\
T265728 & thingi10k & 49 & 0.648 & 0.571 & 0.588 & 0.588 & 0.588 & 0.635 \\
T39108 & thingi10k & 128 & 0.800 & 0.420 & 0.448 & 0.420 & 0.250 & 0.780 \\
T39922 & thingi10k & 168 & 0.667 & 0.964 & 0.667 & 0.667 & 0.661 & 0.667 \\
T40447 & thingi10k & 452 & 1.000 & 1.000 & 1.000 & 1.000 & 1.000 & 1.000 \\
T42025 & thingi10k & 84 & 0.024 & 0.023 & 0.023 & 0.035 & 0.023 & 0.023 \\
T42035 & thingi10k & 154 & 0.987 & 0.808 & 0.957 & 0.727 & 0.807 & 0.662 \\
T48124 & thingi10k & 37 & 0.295 & 0.295 & 0.311 & 0.295 & 0.311 & 0.311 \\
T53749 & thingi10k & 278 & 0.458 & 0.439 & 0.449 & 0.432 & 0.387 & 0.449 \\
T54820 & thingi10k & 82 & 0.427 & 0.427 & 0.427 & 0.546 & 0.427 & 0.427 \\
T56630 & thingi10k & 359 & 0.955 & 0.978 & 0.978 & 0.975 & 0.978 & 0.943 \\
T567026 & thingi10k & 20 & 0.200 & 0.000 & 0.200 & 0.200 & 0.200 & 0.200 \\
T61192 & thingi10k & 26 & 1.000 & 1.000 & 1.000 & 1.000 & 1.000 & 1.000 \\
T69537 & thingi10k & 64 & 0.812 & 0.778 & 0.788 & 0.866 & 0.778 & 0.778 \\
T70559 & thingi10k & 13 & 0.231 & 0.231 & 0.231 & 0.231 & 0.231 & 0.200 \\
T71307 & thingi10k & 76 & 0.237 & 0.224 & 0.272 & 0.350 & 0.350 & 0.330 \\
T75090 & thingi10k & 238 & 0.790 & 0.701 & 0.793 & 0.810 & 0.778 & 0.800 \\
T75651 & thingi10k & 125 & 0.920 & 0.921 & 0.931 & 0.936 & 0.929 & 0.929 \\
T80597 & thingi10k & 188 & 0.883 & 0.819 & 0.851 & 0.856 & 0.841 & 0.815 \\
T81589 & thingi10k & 212 & 0.972 & 0.972 & 0.972 & 0.967 & 0.963 & 0.968 \\
T87521 & thingi10k & 25 & 1.000 & 0.833 & 0.833 & 1.000 & 0.926 & 0.926 \\
T91350 & thingi10k & 433 & 0.993 & 0.984 & 0.986 & 0.989 & 0.973 & 0.984 \\
T91452 & thingi10k & 42 & 0.159 & 0.159 & 0.159 & 0.109 & 0.159 & 0.159 \\
T93881 & thingi10k & 91 & 1.000 & 1.000 & 1.000 & 1.000 & 1.000 & 1.000 \\
T966906 & thingi10k & 265 & 0.817 & 0.814 & 0.879 & 0.773 & 0.896 & 0.782 \\
\midrule
mean & & & 0.705 & 0.672 & 0.693 & 0.693 & 0.686 & 0.690 \\
\bottomrule
\end{tabular}
\end{center}

